\section{Related Literature}\label{sec:lit}

This paper extends the fuzzy DID estimands of \cite{fuzzy2018} by deriving closed-form first-step influence functions for the Wald-DID and time-corrected estimands. The result is a debiased machine-learning formulation with analytic Riesz representers and a data-driven one-sided trimming rule. The contribution sits at the junction of three literatures: debiased ML for sharp DID, IV combined with DID, and influence-function tools for semiparametric efficiency.

In sharp DID, \cite{sant2020doubly} derive the semiparametric efficiency bound for the ATT and propose doubly robust estimators. \cite{chang2020double} bring this score into the DML framework for two-period DID with high-dimensional covariates, and \cite{callaway2023ddml} extend the same logic to staggered adoption with group-time parameters, building on \cite{callaway2021did}. More recent work sharpens efficiency and conditional identification: \cite{chen2025efficient} derive efficient event-study estimators, and \cite{caetano2024did} formalize DID identification under conditional parallel trends. All of these papers study sharp designs, where treatment status coincides with the group-time interaction.

A separate strand combines IV with DID under an external instrument. \cite{ye2023instrumented} develop instrumented DID and propose multiply robust semiparametric estimators for the ATE and CATE, \cite{zhao2025semiparametric} extend the framework to optimal policy learning, and \cite{miyaji2024instrumented} study heterogeneous effects and staggered instrument adoption. These papers do not use debiased machine learning or cross-fitting for high-dimensional covariates. The fuzzy DID design of \cite{fuzzy2018} is different: no external instrument is available, group-time variation itself supplies the instrument, and the target is a LATE for treatment-group switchers rather than an ATE.

Adapting debiased machine learning to the fuzzy design requires a first-step influence function that corrects for ML-estimated nuisances inside a ratio of moments. \cite{hahn1998} provides the classical semiparametric-efficiency benchmark for ATE-type problems, while \cite{newey1994asymptotic} and \cite{ichimura2022influence} give the pathwise-derivative tools behind the FSIF construction. \cite{chernozhukov2022locally} place these ideas in a locally robust GMM framework with cross-fitting, and \cite{CNS2022} and \cite{escanciano2023automatic} estimate Riesz representers numerically through penalized dictionary regressions. This paper instead derives the representers analytically. In the fuzzy DID setting they reduce to propensity-weighted cell indicators, which makes the sign structure transparent and avoids an auxiliary tuning step.
